\section{Stage 1 Brainstorming Notes}\label{app:brainstorming-notes}

During the initial stage of the project, our six-member team conducted a structured ideation session to clarify the problem space surrounding Generative AI in digital content creation. Two brainstorming frameworks were used: the \textbf{\ac{5w1h} method} for problem framing and the \textbf{\ac{scamper} technique} for early-stage solution exploration. The notes below summarize the main directions generated during the collaborative discussion.

\subsection{Problem Framing via 5W1H Matrix}\label{app:5w1h}

\begin{table}[H]
    \centering
    \caption{5W1H Brainstorming Matrix for the Generative AI Content-Creation Dilemma}\label{tab:5w1h}
    \begin{tabularx}{\textwidth}{@{}P{2.2cm}L@{}}
        \toprule
        \textbf{Dimension} & \textbf{Brainstorming Notes and Key Ideas} \\
        \midrule
        \textbf{Who} & AI companies and model developers, data engineers, independent artists, content creators, digital platforms, vulnerable individuals affected by deepfakes, and lawmakers/regulators. \\
        \addlinespace
        \textbf{What} & A growing conflict between large-scale web data collection for AI training, intellectual property protection, and the generation of highly realistic synthetic media that can be beneficial or harmful depending on use. \\
        \addlinespace
        \textbf{When} & The present phase of rapid \ac{genai} commercialization, in which technical capabilities are evolving faster than legal frameworks, policy adaptation, and detection tools. \\
        \addlinespace
        \textbf{Where} & The global digital ecosystem, especially social media platforms, digital art marketplaces, online content-sharing environments, and news-related information spaces. \\
        \addlinespace
        \textbf{Why} & AI organizations require large datasets to improve model capability and maintain competitiveness; creators seek protection of authorship, compensation, and recognition; society requires trustworthy digital information and protection from misuse. \\
        \addlinespace
        \textbf{How} & How can \ac{genai} systems be developed and governed to balance innovation, creator rights, and public trust? Early directions considered included stricter licensing models, provenance tracking, watermarking, platform-level safeguards, and access control mechanisms. \\
        \bottomrule
    \end{tabularx}
\end{table}

\subsection{Early Solution Ideation via SCAMPER}\label{app:scamper}
After the problem was framed, the team used the \ac{scamper} technique to explore possible engineering and governance directions. These ideas were not treated as final solutions, but as exploratory options for later evaluation.

\begin{itemize}
    \item \textbf{Substitute:} Can large-scale scraped copyrighted data be partially substituted with licensed, public-domain, or responsibly curated datasets? \textit{Observation: This may improve legal and ethical compliance, but could reduce dataset diversity and model performance.}

    \item \textbf{Combine:} Can AI data pipelines be combined with royalty-tracking or smart-contract mechanisms? \textit{Idea: A compensation layer could be added for opt-in creators whose works are used in training or referenced in commercial ecosystems.}

    \item \textbf{Adapt:} Can traditional attribution or watermarking practices be adapted for synthetic media? \textit{Idea: Invisible or cryptographic watermarking could help improve provenance and traceability of AI-generated outputs.}

    \item \textbf{Modify:} Can the design of \ac{genai} systems be modified to include stronger safety or provenance controls by default? \textit{Idea: Generation systems could embed output signatures or metadata to support accountability and downstream verification.}

    \item \textbf{Put to another use:} Can the same generative technologies also be used defensively? \textit{Idea: AI could support detection, moderation, or authenticity verification systems rather than only content generation.}

    \item \textbf{Eliminate:} What happens if copyrighted or uncertain-source data is eliminated entirely from training pipelines? \textit{Observation: This may increase compliance and reduce disputes, but could also shrink training data availability and affect model capability.}

    \item \textbf{Reverse / Rearrange:} Instead of assuming unrestricted model access first and adding safeguards later, can governance be reversed so that provenance, consent, and accountability are built into the pipeline from the beginning? \textit{Idea: A ``compliance-by-design'' approach may be more sustainable than retroactive regulation.}
\end{itemize}

\subsection{Brainstorming Outcome}\label{app:brainstorming-outcome}
The brainstorming session helped the team move from a broad discussion about ``AI versus human creators'' toward a more structured understanding of the issue as a \textbf{socio-technical engineering challenge}. In particular, the \ac{5w1h} analysis clarified the affected stakeholders, the urgency of the problem, and the competing interests involved, while \ac{scamper} helped identify possible intervention points in data governance, content provenance, platform responsibility, and system design. These notes informed the more formal problem framing in \cref{sec:brainstorming} of the main report.
